\documentclass{article}

% Language setting
\usepackage[english]{babel}

% Page size and margins
\usepackage[letterpaper,top=2cm,bottom=2cm,left=3cm,right=3cm,marginparwidth=1.75cm]{geometry}

% Math and theorem packages
\usepackage{amsmath}
\usepackage{amssymb}
\usepackage{amsthm}

% Bibliography
\usepackage{biblatex}
\addbibresource{reference.bib}

% Hyperlinks (load last)
\usepackage[colorlinks=true, allcolors=blue]{hyperref}

% Theorem environments
\theoremstyle{plain}
\newtheorem{theorem}{Theorem}[section]
\theoremstyle{definition}
\newtheorem{defi}{Definition}[section]
\theoremstyle{remark}
\newtheorem*{remark}{Remark}

% Custom operators
\DeclareMathOperator*{\argmax}{arg\,max}
\DeclareMathOperator{\val}{val}
\DeclareMathOperator{\nextplayer}{next}

\title{Modeling the Game Flip 7}
\author{Jesse Georgias}

\begin{document}
\maketitle

\begin{abstract}
This paper presents a game-theoretic model for the card game ``Flip 7,'' which was obtained by Bruno in a McDonald's Happy Meal. Starting from a single-player Markov Decision Process for a simplified deck, we incrementally incorporate sequential two-player competition, bonus modifier cards, action cards, and finally the full $n$-player round-robin game modeled as a finite extensive-form game solved by backward induction.
% Replace the URL below with the official rules link if submitting formally.
The game rules can be read \href{https://www.youtube.com/watch?v=dQw4w9WgXcQ}{here}.
\end{abstract}

\section{The Framework}

\subsection{Markov Decision Processes}
We begin by modeling the single-player decision problem via Markov Decision Processes (MDPs) and solving the resulting expected-payoff optimization via the Bellman Optimality Equation.

\begin{defi}[\cite{sutton2018reinforcement}]
A Markov Decision Process (MDP) is formally defined as a 5-tuple
\[
    \mathcal{M} = (\mathcal{S}, \mathcal{A}, \mathcal{P}, \mathcal{R}, \gamma),
\]
where:
\begin{enumerate}
    \item $\mathcal{S}$ is the set of valid states (the state space).
    \item $\mathcal{A}$ is the set of valid actions (the action space).
    \item $\mathcal{P}: \mathcal{S} \times \mathcal{A} \times \mathcal{S} \to [0, 1]$ is the state transition probability function, defining the probability of transitioning to state $s'$ given state $s$ and action $a$:
    \[
        \mathcal{P}(s' \mid s, a) = \Pr(S_{t+1} = s' \mid S_t = s, A_t = a),
    \]
    subject to $\sum_{s' \in \mathcal{S}} \mathcal{P}(s' \mid s, a) = 1$ for all $(s, a) \in \mathcal{S} \times \mathcal{A}$.
    \item $\mathcal{R}: \mathcal{S} \times \mathcal{A} \times \mathcal{S} \to \mathbb{R}$ is the reward function, representing the expected immediate reward received after transitioning from state $s$ to state $s'$ via action $a$.
    \item $\gamma \in [0, 1]$ is the discount factor, which determines the present value of future rewards.
\end{enumerate}
\end{defi}

The domain of the transition function $\mathcal{P}$ and the reward function $\mathcal{R}$ inherently enforces the \textit{Markov property}: future dynamics depend solely on the current state and action, rendering the next state conditionally independent of the prior history:
\[
    \Pr(S_{t+1} = s' \mid S_t = s, A_t = a, S_{t-1}, A_{t-1}, \dots, S_0, A_0) = \Pr(S_{t+1} = s' \mid S_t = s, A_t = a).
\]

The goal in solving an MDP is to find an optimal \textit{policy} $\pi$ (for deterministic policies, a mapping $\pi: \mathcal{S} \to \mathcal{A}$) that maximizes the expected discounted return $G_t = \sum_{k=0}^{\infty} \gamma^k R_{t+k+1}$ from every state, i.e., maximizes $\mathbb{E}_{\pi}[G_t \mid S_t = s]$ for all $s \in \mathcal{S}$.

\begin{theorem}[Bellman Optimality Theorem]\label{thm:bellman}
For any finite Markov Decision Process with a discount factor $\gamma \in [0, 1)$, there exists a unique optimal state-value function $V^* : \mathcal{S} \to \mathbb{R}$ satisfying the Bellman Optimality Equation
\[
    V^*(s) = \max_{a \in \mathcal{A}} \sum_{s' \in \mathcal{S}} \mathcal{P}(s' \mid s, a) \left[ \mathcal{R}(s, a, s') + \gamma V^*(s') \right]
\]
for all $s \in \mathcal{S}$. Furthermore, any policy $\pi^*$ that selects actions greedily with respect to $V^*$ is an optimal policy.
\begin{proof}
See Theorem~6.2.3 of \cite{puterman2014markov} or Chapter~4 of \cite{sutton2018reinforcement}.
\end{proof}
\end{theorem}

\begin{remark}[Undiscounted Finite Horizons ($\gamma = 1$)]
While Theorem~\ref{thm:bellman} requires $\gamma \in [0, 1)$ to guarantee contraction in infinite-horizon MDPs, a single round of Flip 7 without mid-round reshuffling is a finite-horizon, directed acyclic MDP: every $\text{draw}$ action strictly decreases the finite deck size $|D|$, and every $\text{stay}$ action transitions immediately to an absorbing terminal state. Consequently, the existence and uniqueness of $V^*$ remain guaranteed for $\gamma = 1$ via finite backward induction on $|D|$.
\end{remark}

\subsection{The Simplified Single-Player Game}
A standard Flip 7 deck consists of $94$ cards:
\begin{itemize}
    \item $79$ \textbf{number cards} from $\mathcal{N} = \{0, 1, \dots, 12\}$, containing one $0$ card and $x$ copies of each integer $x \in \{1, \dots, 12\}$ ($1 + \sum_{x=1}^{12} x = 79$);
    \item $6$ \textbf{bonus modifier cards}, consisting of one multiplier card ($\times 2$) and five flat bonus cards ($+2, +4, +6, +8, +10$);
    \item $9$ \textbf{action cards}, consisting of three copies each of \textit{Freeze}, \textit{Second Chance}, and \textit{Flip Three}.
\end{itemize}
Throughout this paper, we analyze a single round of play under the assumption that no mid-round reshuffle occurs; if the draw pile becomes empty ($|D| = 0$), no further draws are possible and active players retain their current tableaus.

To build our baseline model, we first exclude all bonus and action cards and solve the turn of a single player as an MDP $\mathcal{M} = (\mathcal{S}, \mathcal{A}, \mathcal{P}, \mathcal{R}, \gamma)$:

\begin{enumerate}
    \item \textbf{State Space $\mathcal{S}$:} Let $\mathcal{S} = \big(2^{\mathcal{N}} \times \mathbb{N}_0^{13}\big) \cup \{s_{\text{term}}\}$, where $s_{\text{term}}$ is an absorbing terminal state denoting the end of the turn, and every non-terminal state is a tuple $s = (C, D)$ with:
    \begin{itemize}
        \item $C \subseteq \mathcal{N}$: the set of unique number cards currently in front of the player (with $|C| \le 6$ in any non-terminal state);
        \item $D = (d_0, d_1, \dots, d_{12}) \in \mathbb{N}_0^{13}$: the count vector of cards remaining in the deck.
    \end{itemize}
    The total number of remaining cards is $|D| = \sum_{x=0}^{12} d_x$. Let $\boldsymbol{\delta}_x \in \mathbb{N}_0^{13}$ denote the unit basis vector with a $1$ in entry $x$ and $0$ elsewhere.

    \item \textbf{Action Space $\mathcal{A}$:} For non-terminal states with $|D| > 0$, the available action set is binary:
    \[
        \mathcal{A} = \{\text{stay}, \text{draw}\} \cong \{0, 1\}.
    \]
    When $|D| = 0$, only $a = \text{stay}$ is available.

    \item \textbf{Transition Probabilities $\mathcal{P}$:} From a non-terminal state $s = (C, D)$:
    \begin{itemize}
        \item If $a = \text{stay}$, the turn terminates deterministically: $\mathcal{P}(s_{\text{term}} \mid s, \text{stay}) = 1$.
        \item If $a = \text{draw}$ (requiring $|D| > 0$), a card $x \in \mathcal{N}$ is drawn with probability $d_x / |D|$, updating the deck to $D' = D - \boldsymbol{\delta}_x$. If $x \in C$ (bust) or $x \notin C$ with $|C| = 6$ (Flip 7), the process transitions to $s_{\text{term}}$; otherwise it transitions to $s' = (C \cup \{x\}, D - \boldsymbol{\delta}_x)$.
    \end{itemize}
    The terminal state is absorbing: $\mathcal{P}(s_{\text{term}} \mid s_{\text{term}}, a) = 1$ with zero subsequent reward.

    \item \textbf{Reward Function $\mathcal{R}$:} For a non-terminal state $s = (C, D)$, immediate rewards are awarded only upon termination:
    \begin{itemize}
        \item \textit{Stopping Safely ($a = \text{stay}$):} $\mathcal{R}(s, \text{stay}, s_{\text{term}}) = \sum_{y \in C} y$ (with the empty sum equal to $0$).
        \item \textit{Busting ($a = \text{draw},\, x \in C$):} the turn ends with reward $0$.
        \item \textit{Flip 7 ($a = \text{draw},\, x \notin C,\, |C| = 6$):} the turn ends immediately with reward $\sum_{y \in C} y + x + 15$.
        \item \textit{Safe Draw ($a = \text{draw},\, x \notin C,\, |C| < 6$):} the turn continues to $(C \cup \{x\}, D - \boldsymbol{\delta}_x)$ with immediate reward $0$.
    \end{itemize}

    \item \textbf{Discount Factor $\gamma$:} Because the turn is guaranteed to terminate in at most $\min(7, |D| + 1)$ steps, we set $\gamma = 1$.
\end{enumerate}

\subsection{Bellman Optimum for the Simplified Game}
For notational clarity, we drop the asterisk on $V^*$ henceforth and write $V(C, D)$ for the optimal expected score from non-terminal state $(C, D)$. When the deck is exhausted ($|D| = 0$), the player must stop and secures $V(C, \mathbf{0}) = \sum_{y \in C} y$. For $|D| > 0$, the Bellman optimality equation is
\[
    V(C, D) = \max \Big( Q_{\text{stay}}(C, D), \, Q_{\text{draw}}(C, D) \Big),
\]
where
\[
    Q_{\text{stay}}(C, D) = \sum_{y \in C} y, \qquad
    Q_{\text{draw}}(C, D) = \sum_{\substack{x \in \mathcal{N} \\ d_x > 0}} \frac{d_x}{|D|} \cdot v_{\text{next}}(C, D, x),
\]
and the post-draw continuation value $v_{\text{next}}(C, D, x)$ encodes busting, achieving a Flip 7, or continuing safely:
\[
    v_{\text{next}}(C, D, x) =
    \begin{cases}
        0 & \text{if } x \in C \text{ (Bust)}, \\
        \sum_{y \in C} y + x + 15 & \text{if } x \notin C \text{ and } |C| = 6 \text{ (Flip 7)}, \\
        V(C \cup \{x\}, D - \boldsymbol{\delta}_x) & \text{if } x \notin C \text{ and } |C| < 6 \text{ (Safe Draw)}.
    \end{cases}
\]

\section{The Two-Player Sequential Game}
Before modeling the true round-robin single-card draw structure of Flip 7 in Section~\ref{sec:nplayer}, we develop a tractable two-player \textit{sequential baseline} in which Player~1 completes their entire turn before Player~2 draws. Because drawn cards are removed from the draw pile, Player~2 inherits a depleted deck shaped by Player~1's turn. Assuming both players seek to maximize their expected point differential ($S_1 - S_2$ for Player~1 and $S_2 - S_1$ for Player~2), the resulting two-player game is strictly zero-sum and can be solved via backward induction: we first compute Player~2's optimal value function across all remaining deck states, then fold those values into Player~1's Bellman equation.

\subsection{The Simplified Two-Player Game}
We enumerate the players as $\{1, 2\}$ and define value functions $V_1, V_2: 2^{\mathcal{N}} \times \mathbb{N}_0^{13} \to \mathbb{R}$.

\begin{enumerate}
    \item \textbf{Player~2's Value Function ($V_2$):} Once Player~1 has finished their turn, Player~1's realized score $S_1$ is a fixed constant. Maximizing expected differential $\mathbb{E}[S_2 - S_1] = \mathbb{E}[S_2] - S_1$ is therefore equivalent to maximizing Player~2's own expected score $\mathbb{E}[S_2]$. Thus $V_2(C, D)$ is identical to the single-player value function $V(C, D)$ of Section~1.3, with base case $V_2(C, \mathbf{0}) = \sum_{y \in C} y$ and, for $|D| > 0$,
    \[
        V_2(C, D) = \max \left( \sum_{y \in C} y, \sum_{\substack{x \in \mathcal{N} \\ d_x > 0}} \frac{d_x}{|D|} \cdot v_{2,\text{next}}(C, D, x) \right),
    \]
    where $v_{2,\text{next}}(C, D, x)$ matches $v_{\text{next}}(C, D, x)$ from Section~1.3.

    \item \textbf{Player~1's Value Function ($V_1$):} Let $V_1(C, D)$ denote Player~1's optimal expected point differential $S_1 - S_2$ from state $(C, D)$. If $|D| = 0$, Player~1 must stop and Player~2 cannot draw, yielding $V_1(C, \mathbf{0}) = \sum_{y \in C} y$. For $|D| > 0$,
    \[
        V_1(C, D) = \max \Big( Q_{1,\text{stay}}(C, D), \, Q_{1,\text{draw}}(C, D) \Big),
    \]
    with:
    \begin{itemize}
        \item \textbf{Stay:} Player~1 secures $\sum_{y \in C} y$, and Player~2 begins with an empty tableau ($\emptyset$) and remaining deck $D$:
        \[
            Q_{1,\text{stay}}(C, D) = \sum_{y \in C} y - V_2(\emptyset, D).
        \]
        \item \textbf{Draw:}
        \[
            Q_{1,\text{draw}}(C, D) = \sum_{\substack{x \in \mathcal{N} \\ d_x > 0}} \frac{d_x}{|D|} \cdot v_{1,\text{next}}(C, D, x),
        \]
        where a Flip 7 ends the round immediately (leaving Player~2 with $0$ points), while a bust leaves Player~2 to play on $D - \boldsymbol{\delta}_x$:
        \[
            v_{1,\text{next}}(C, D, x) =
            \begin{cases}
                0 - V_2(\emptyset, D - \boldsymbol{\delta}_x) & \text{if } x \in C \text{ (Bust)}, \\
                \sum_{y \in C} y + x + 15 & \text{if } x \notin C \text{ and } |C| = 6 \text{ (Flip 7)}, \\
                V_1(C \cup \{x\}, D - \boldsymbol{\delta}_x) & \text{if } x \notin C \text{ and } |C| < 6 \text{ (Safe Draw)}.
            \end{cases}
        \]
    \end{itemize}
\end{enumerate}

\subsection{Incorporating Bonus Cards}
We next introduce the $6$ bonus modifier cards while still excluding action cards. Bonus cards alter a player's final score but can never cause a bust and do not count toward the $7$ unique number cards required for a Flip 7.

\subsubsection{Expanded State and Action Space}
\begin{enumerate}
    \item \textbf{Card Types:} Let $\mathcal{B} = \mathcal{B}_{\times} \cup \mathcal{B}_{+}$ be the set of $6$ bonus card types, partitioned into the single multiplier card $\mathcal{B}_{\times} = \{\text{x2}\}$ and the five flat bonus cards $\mathcal{B}_{+} = \{\text{+2}, \text{+4}, \text{+6}, \text{+8}, \text{+10}\}$, equipped with the point-value map $\val: \mathcal{B}_{+} \to \{2, 4, 6, 8, 10\}$. The set of drawable card types is $X = \mathcal{N} \cup \mathcal{B}$.

    \item \textbf{State Space:} A non-terminal state for the current player is a 4-tuple $s = (C, m, p, D)$, where:
    \begin{itemize}
        \item $C \subseteq \mathcal{N}$: the set of unique number cards in front of the player;
        \item $m \in \{1, 2\}$: the active number-card multiplier (initialized to $1$, becoming $2$ when the unique $\text{x2}$ card is drawn);
        \item $p \in \mathbb{N}_0$: the sum of accumulated flat bonus points;
        \item $D = (d_x)_{x \in X} \in \mathbb{N}_0^{|X|}$: the remaining card counts in the deck, with $|D| = \sum_{x \in X} d_x$ and unit basis vectors $\boldsymbol{\delta}_x \in \mathbb{N}_0^{|X|}$.
    \end{itemize}

    \item \textbf{Action Space:} $\mathcal{A} = \{\text{stay}, \text{draw}\}$ for $|D| > 0$, and $\{\text{stay}\}$ for $|D| = 0$.
\end{enumerate}

\subsubsection{Reward Function and Scoring}
Scoring obeys a strict order of operations: (1) sum the number cards, (2) apply the multiplier $m$, (3) add flat bonus points $p$, and (4) add $15$ points if a Flip 7 is achieved. Thus, the secured score upon stopping safely is
\[
    R_{\text{safe}}(C, m, p) = m \sum_{y \in C} y + p,
\]
and the score upon achieving a Flip 7 by drawing a 7th unique number card $x \in \mathcal{N} \setminus C$ (when $|C| = 6$) is
\[
    R_{\text{flip7}}(C, m, p, x) = m \left( \sum_{y \in C} y + x \right) + p + 15.
\]

\subsubsection{Bellman Optimum for the Single-Player Game with Bonus Cards}
With base case $V(C, m, p, \mathbf{0}) = R_{\text{safe}}(C, m, p)$, the optimal expected score for $|D| > 0$ satisfies
\[
    V(C, m, p, D) = \max \Big( R_{\text{safe}}(C, m, p), \, Q_{\text{draw}}(C, m, p, D) \Big),
\]
where
\[
    Q_{\text{draw}}(C, m, p, D) = \sum_{\substack{x \in X \\ d_x > 0}} \frac{d_x}{|D|} \cdot v_{\text{next}}(C, m, p, D, x),
\]
and
\[
    v_{\text{next}}(C, m, p, D, x) =
    \begin{cases}
        0 & \text{if } x \in \mathcal{N} \text{ and } x \in C \text{ (Bust)}, \\
        R_{\text{flip7}}(C, m, p, x) & \text{if } x \in \mathcal{N},\, x \notin C,\text{ and } |C| = 6 \text{ (Flip 7)}, \\
        V(C \cup \{x\}, m, p, D - \boldsymbol{\delta}_x) & \text{if } x \in \mathcal{N},\, x \notin C,\text{ and } |C| < 6 \text{ (Safe Draw)}, \\
        V(C, 2m, p, D - \boldsymbol{\delta}_x) & \text{if } x \in \mathcal{B}_{\times} \text{ (Multiplier Draw)}, \\
        V(C, m, p + \val(x), D - \boldsymbol{\delta}_x) & \text{if } x \in \mathcal{B}_{+} \text{ (Flat Bonus Draw)}.
    \end{cases}
\]

\subsubsection{The Two-Player Game with Bonus Cards}
Player~2's optimal expected score $V_2(C, m, p, D)$ coincides with the single-player function $V(C, m, p, D)$ above. For Player~1, the deck-exhaustion base case is $V_1(C, m, p, \mathbf{0}) = R_{\text{safe}}(C, m, p)$. For $|D| > 0$,
\[
    V_1(C, m, p, D) = \max \Big( Q_{1,\text{stay}}(C, m, p, D), \, Q_{1,\text{draw}}(C, m, p, D) \Big),
\]
where
\[
    Q_{1,\text{stay}}(C, m, p, D) = R_{\text{safe}}(C, m, p) - V_2(\emptyset, 1, 0, D),
\]
and $Q_{1,\text{draw}}(C, m, p, D) = \sum_{x \in X : d_x > 0} \frac{d_x}{|D|} \cdot v_{1,\text{next}}(C, m, p, D, x)$, with
\[
    v_{1,\text{next}}(C, m, p, D, x) =
    \begin{cases}
        0 - V_2(\emptyset, 1, 0, D - \boldsymbol{\delta}_x) & \text{if } x \in \mathcal{N},\, x \in C \text{ (Bust)}, \\
        R_{\text{flip7}}(C, m, p, x) & \text{if } x \in \mathcal{N},\, x \notin C,\, |C| = 6 \text{ (Flip 7)}, \\
        V_1(C \cup \{x\}, m, p, D - \boldsymbol{\delta}_x) & \text{if } x \in \mathcal{N},\, x \notin C,\, |C| < 6, \\
        V_1(C, 2m, p, D - \boldsymbol{\delta}_x) & \text{if } x \in \mathcal{B}_{\times}, \\
        V_1(C, m, p + \val(x), D - \boldsymbol{\delta}_x) & \text{if } x \in \mathcal{B}_{+}.
    \end{cases}
\]

\subsection{Incorporating Action Cards: The Joint State Space}
Introducing the action cards $\mathcal{A}_{\text{cards}} = \{\text{fr}, \text{sc}, \text{f3}\}$ (Freeze, Second Chance, and Flip Three) couples the players' states because cards drawn during Player~1's turn can immediately alter Player~2's tableau or force Player~2 to draw.

\subsubsection{The Joint State Space and Shorthand Notation}
Let the complete set of card types be $X = \mathcal{N} \cup \mathcal{B} \cup \mathcal{A}_{\text{cards}}$. A joint state is a tuple $s = (E_1, E_2, D)$, where:
\begin{enumerate}
    \item \textbf{Deck Vector $D$:} $D = (d_x)_{x \in X} \in \mathbb{N}_0^{|X|}$ tracks the remaining count of each of the $22$ distinct card types ($13$ number, $6$ bonus, and $3$ action types, initially summing to $94$ cards), with $|D| = \sum_{x \in X} d_x$ and unit basis vectors $\boldsymbol{\delta}_x \in \mathbb{N}_0^{|X|}$.
    \item \textbf{Player Environment $E_i$:} For each player $i \in \{1, 2\}$,
    \[
        E_i = (C_i, m_i, p_i, sc_i, \alpha_i),
    \]
    where $C_i \subseteq \mathcal{N}$, $m_i \in \{1, 2\}$, $p_i \in \mathbb{N}_0$, $sc_i \in \{0, 1\}$ indicates possession of an active Second Chance card, and $\alpha_i \in \{0, 1\}$ indicates whether Player~$i$ is still \textit{active} in the round ($\alpha_i = 1$) or has stopped, busted, or been frozen ($\alpha_i = 0$).
\end{enumerate}
We overload the scoring functions to accept a player environment directly:
\[
    R_{\text{safe}}(E_i) = m_i \sum_{y \in C_i} y + p_i, \qquad
    R_{\text{flip7}}(E_i, x) = m_i \left( \sum_{y \in C_i} y + x \right) + p_i + 15.
\]
When a player busts, their environment resets to $E^{\text{bust}} = (\emptyset, 1, 0, 0, 0)$, which satisfies $\alpha_i = 0$ and $R_{\text{safe}}(E^{\text{bust}}) = 0$. We write $E_i[\alpha_i \leftarrow 0]$ (and similarly for other fields) to denote $E_i$ with the specified coordinate replaced.

\subsubsection{Opponent-First Targeting Assumption}
Under official Flip 7 rules, any active player including oneself may be targeted by a Freeze or Flip Three card, and self-targeting can occasionally be optimal (e.g., freezing oneself to avoid a risky draw or avoiding giving an opponent with six cards extra chances at a Flip 7). To keep the two-player sequential baseline compact, we adopt an \textit{opponent-first targeting assumption}: whenever an active opponent exists, Freeze and Flip Three automatically target the opponent; only when the acting player is the sole remaining active player do they target themselves. (This assumption is relaxed easily by replacing the fixed target with a choice over $\{1, 2\}$ whenever both players are active.)

\subsubsection{Player~2's Solo Value Function with Action Cards}
When Player~1's turn concludes (whether by staying, freezing, or busting), Player~1 is inactive ($\alpha_1 = 0$), so Player~2 acts as the sole active player. If $\alpha_2 = 0$ (Player~2 was already frozen or busted during Player~1's turn) or $|D| = 0$, Player~2 cannot draw and secures $V_2(E_2, D) = R_{\text{safe}}(E_2)$. For $\alpha_2 = 1$ and $|D| > 0$,
\[
    V_2(E_2, D) = \max \left( R_{\text{safe}}(E_2), \sum_{\substack{x \in X \\ d_x > 0}} \frac{d_x}{|D|} \cdot v_{2,\text{next}}(E_2, D, x) \right),
\]
where:
\begin{itemize}
    \item \textbf{Duplicate Number ($x \in \mathcal{N} \cap C_2$):} $V_2(E_2[sc_2 \leftarrow 0], D - \boldsymbol{\delta}_x)$ if $sc_2 = 1$; otherwise $0$ (bust).
    \item \textbf{Safe Number ($x \in \mathcal{N} \setminus C_2$):} $R_{\text{flip7}}(E_2, x)$ if $|C_2| = 6$; otherwise $V_2(E_2[C_2 \leftarrow C_2 \cup \{x\}], D - \boldsymbol{\delta}_x)$.
    \item \textbf{Bonus Card ($x \in \mathcal{B}$):} $V_2(E_2[m_2 \leftarrow 2m_2], D - \boldsymbol{\delta}_x)$ if $x \in \mathcal{B}_{\times}$, and $V_2(E_2[p_2 \leftarrow p_2 + \val(x)], D - \boldsymbol{\delta}_x)$ if $x \in \mathcal{B}_{+}$.
    \item \textbf{Second Chance ($x = \text{sc}$):} $V_2(E_2[sc_2 \leftarrow 1], D - \boldsymbol{\delta}_x)$ (a surplus copy when $sc_2 = 1$ is discarded since no other active player exists).
    \item \textbf{Freeze ($x = \text{fr}$):} Player~2 must freeze themselves, immediately securing $R_{\text{safe}}(E_2)$.
    \item \textbf{Flip Three ($x = \text{f3}$):} Player~2 must undergo $3$ forced draws via the solo recursion $W_3^{\text{solo}}(E_2, D - \boldsymbol{\delta}_x, \emptyset; V_2)$ defined in Section~2.4.
\end{itemize}

\subsubsection{Player~1's Bellman Equation with Action Cards}
Player~1 maximizes the expected point differential $V_1(E_1, E_2, D)$. If Player~1 is inactive ($\alpha_1 = 0$, e.g., frozen or busted during a Flip Three resolution) or the deck is exhausted ($|D| = 0$), Player~1's turn ends immediately and Player~2 completes the round:
\[
    V_1(E_1, E_2, D) = R_{\text{safe}}(E_1) - V_2(E_2, D) \qquad \text{if } \alpha_1 = 0 \text{ or } |D| = 0.
\]
When $\alpha_1 = 1$ and $|D| > 0$,
\[
    V_1(E_1, E_2, D) = \max \Big( Q_{1,\text{stay}}(E_1, E_2, D), \, Q_{1,\text{draw}}(E_1, E_2, D) \Big),
\]
where
\[
    Q_{1,\text{stay}}(E_1, E_2, D) = R_{\text{safe}}(E_1) - V_2(E_2, D), \qquad
    Q_{1,\text{draw}}(E_1, E_2, D) = \sum_{\substack{x \in X \\ d_x > 0}} \frac{d_x}{|D|} \cdot v_{1,\text{next}}(E_1, E_2, D, x).
\]
The post-draw continuation $v_{1,\text{next}}(E_1, E_2, D, x)$ evaluates each card type as follows:
\begin{itemize}
    \item \textbf{Duplicate Number Card ($x \in \mathcal{N} \cap C_1$):}
    \[
        v_{1,\text{next}} =
        \begin{cases}
            V_1\big(E_1[sc_1 \leftarrow 0], E_2, D - \boldsymbol{\delta}_x\big) & \text{if } sc_1 = 1 \text{ (Protected)}, \\
            0 - V_2(E_2, D - \boldsymbol{\delta}_x) & \text{if } sc_1 = 0 \text{ (Bust)}.
        \end{cases}
    \]

    \item \textbf{Safe Number Card ($x \in \mathcal{N} \setminus C_1$):}
    \[
        v_{1,\text{next}} =
        \begin{cases}
            R_{\text{flip7}}(E_1, x) - R_{\text{safe}}(E_2) & \text{if } |C_1| = 6 \text{ (Flip 7)}, \\
            V_1\big(E_1[C_1 \leftarrow C_1 \cup \{x\}], E_2, D - \boldsymbol{\delta}_x\big) & \text{if } |C_1| < 6 \text{ (Safe Draw)}.
        \end{cases}
    \]

    \item \textbf{Bonus Card ($x \in \mathcal{B}$):}
    \[
        v_{1,\text{next}} =
        \begin{cases}
            V_1\big(E_1[m_1 \leftarrow 2m_1], E_2, D - \boldsymbol{\delta}_x\big) & \text{if } x \in \mathcal{B}_{\times}, \\
            V_1\big(E_1[p_1 \leftarrow p_1 + \val(x)], E_2, D - \boldsymbol{\delta}_x\big) & \text{if } x \in \mathcal{B}_{+}.
        \end{cases}
    \]

    \item \textbf{Second Chance Card ($x = \text{sc}$):} A player may hold at most one active Second Chance card; a duplicate is passed to an active opponent who does not already hold one, or discarded otherwise:
    \[
        v_{1,\text{next}} =
        \begin{cases}
            V_1\big(E_1[sc_1 \leftarrow 1], E_2, D - \boldsymbol{\delta}_x\big) & \text{if } sc_1 = 0, \\
            V_1\big(E_1, E_2[sc_2 \leftarrow 1], D - \boldsymbol{\delta}_x\big) & \text{if } sc_1 = 1,\, \alpha_2 = 1,\, sc_2 = 0, \\
            V_1\big(E_1, E_2, D - \boldsymbol{\delta}_x\big) & \text{if } sc_1 = 1 \text{ and } (\alpha_2 = 0 \text{ or } sc_2 = 1).
        \end{cases}
    \]

    \item \textbf{Freeze Card ($x = \text{fr}$):} Targets Player~2 if active ($\alpha_2 = 1$), or freezes Player~1 if Player~2 is already inactive ($\alpha_2 = 0$):
    \[
        v_{1,\text{next}} =
        \begin{cases}
            V_1\big(E_1, E_2[\alpha_2 \leftarrow 0], D - \boldsymbol{\delta}_x\big) & \text{if } \alpha_2 = 1, \\
            R_{\text{safe}}(E_1) - R_{\text{safe}}(E_2) & \text{if } \alpha_2 = 0.
        \end{cases}
    \]

    \item \textbf{Flip Three Card ($x = \text{f3}$):} Forces Player~2 to draw $3$ cards if $\alpha_2 = 1$, or forces Player~1 to draw $3$ cards if $\alpha_2 = 0$:
    \[
        v_{1,\text{next}} =
        \begin{cases}
            W_3^{(2)}\big(E_1, E_2, D - \boldsymbol{\delta}_x, \emptyset; \, V_1\big) & \text{if } \alpha_2 = 1, \\
            W_3^{(1)}\big(E_1, E_2, D - \boldsymbol{\delta}_x, \emptyset; \, V_1\big) & \text{if } \alpha_2 = 0.
        \end{cases}
    \]
\end{itemize}

\subsection{Expanding the Flip Three Sub-Game}
When a player $t \in \{1, 2\}$ is targeted by a Flip Three card, they must draw up to $3$ cards consecutively. In our modeled ruleset, this introduces two mechanics:
\begin{enumerate}
    \item \textbf{Delayed Execution:} Any action cards drawn during the $3$ forced draws are collected into a temporary multiset $A_{\text{temp}}$ and resolved only after the draws finish.
    \item \textbf{Action-Card Transfer on Bust:} If target $t$ busts during the forced draws, they score $0$, become inactive ($\alpha_t \leftarrow 0$), and transfer $A_{\text{temp}}$ to the other player $-t = 3 - t$ (if active).
\end{enumerate}
\begin{remark}
Some printings or house rules of Flip 7 activate a drawn Second Chance immediately during Flip Three or discard unresolved Freeze/Flip Three cards if the target busts mid-draw. Both variants are obtained directly from the equations below by moving the $x = \text{sc}$ branch out of $A_{\text{temp}}$ or replacing the bust transfer with a direct call to the continuation $\kappa$.
\end{remark}

Because resolving $A_{\text{temp}}$ can trigger nested Flip Three sub-games, we parameterize both the forced-draw recursion $W_k^{(t)}(E_1, E_2, D, A_{\text{temp}}; \kappa)$ and the resolution operator $\Psi_r(E_1, E_2, D, A_{\text{temp}}; \kappa)$ by a \textit{continuation function} $\kappa(E_1, E_2, D)$ specifying the expected differential once the current sub-game finishes. At the top level of Player~1's turn, $\kappa = V_1$.

\subsubsection{The Recursive Forced-Draw Step ($W_k^{(t)}$)}
If $k = 0$ or the deck is exhausted ($|D| = 0$), the forced-draw phase ends and target $t$ resolves $A_{\text{temp}}$:
\[
    W_k^{(t)}(E_1, E_2, D, A_{\text{temp}}; \kappa) = \Psi_t(E_1, E_2, D, A_{\text{temp}}; \kappa) \qquad \text{if } k = 0 \text{ or } |D| = 0.
\]
For $k \in \{1, 2, 3\}$ and $|D| > 0$, target $t$ draws a card $x \in X$ with probability $d_x / |D|$:
\[
    W_k^{(t)}(E_1, E_2, D, A_{\text{temp}}; \kappa) = \sum_{\substack{x \in X \\ d_x > 0}} \frac{d_x}{|D|} \cdot w_{\text{step}}^{(t)}(x),
\]
where, letting $-t = 3 - t$ and writing $(E_1', E_2')$ for the pair updated in coordinate $t$:
\begin{itemize}
    \item \textbf{Action Card ($x \in \mathcal{A}_{\text{cards}}$):} $w_{\text{step}}^{(t)}(x) = W_{k-1}^{(t)}(E_1, E_2, D - \boldsymbol{\delta}_x, A_{\text{temp}} \cup \{x\}; \kappa)$.
    \item \textbf{Bonus Card ($x \in \mathcal{B}$):} update $E_t$ to $E_t[m_t \leftarrow 2m_t]$ (if $x \in \mathcal{B}_{\times}$) or $E_t[p_t \leftarrow p_t + \val(x)]$ (if $x \in \mathcal{B}_{+}$) and set $w_{\text{step}}^{(t)}(x) = W_{k-1}^{(t)}(E_1', E_2', D - \boldsymbol{\delta}_x, A_{\text{temp}}; \kappa)$.
    \item \textbf{Safe Number Card ($x \in \mathcal{N} \setminus C_t$):} if $|C_t| = 6$, the round ends immediately with differential
    \[
        w_{\text{step}}^{(t)}(x) =
        \begin{cases}
            R_{\text{flip7}}(E_1, x) - R_{\text{safe}}(E_2) & \text{if } t = 1, \\
            R_{\text{safe}}(E_1) - R_{\text{flip7}}(E_2, x) & \text{if } t = 2;
        \end{cases}
    \]
    otherwise update $C_t \leftarrow C_t \cup \{x\}$ and set $w_{\text{step}}^{(t)}(x) = W_{k-1}^{(t)}(E_1', E_2', D - \boldsymbol{\delta}_x, A_{\text{temp}}; \kappa)$.
    \item \textbf{Duplicate Number Card ($x \in \mathcal{N} \cap C_t$):}
    \begin{itemize}
        \item If $sc_t = 1$ (Protected): update $sc_t \leftarrow 0$ and set $w_{\text{step}}^{(t)}(x) = W_{k-1}^{(t)}(E_1', E_2', D - \boldsymbol{\delta}_x, A_{\text{temp}}; \kappa)$.
        \item If $sc_t = 0$ (Bust \& Transfer): set $E_t' = E^{\text{bust}} = (\emptyset, 1, 0, 0, 0)$ (leaving $E_{-t}' = E_{-t}$). If $\alpha_{-t} = 1$, transfer $A_{\text{temp}}$ to $-t$ via $\Psi_{-t}(E_1', E_2', D - \boldsymbol{\delta}_x, A_{\text{temp}}; \kappa)$; if $\alpha_{-t} = 0$, no active players remain and the round ends immediately with $R_{\text{safe}}(E_1') - R_{\text{safe}}(E_2')$.
    \end{itemize}
\end{itemize}

\subsubsection{The Action-Card Resolution Operator ($\Psi_r$)}
When multiple action cards accumulate in $A_{\text{temp}}$, their execution order affects the expected outcome (for example, playing Flip Three before Freeze can force a bust, whereas Freezing first locks in the opponent's points). Because the two-player game is zero-sum, Player~1 chooses the next card $c \in A_{\text{temp}}$ to maximize the differential ($\operatorname{opt}_1 = \max$), whereas Player~2 chooses to minimize it ($\operatorname{opt}_2 = \min$).

For resolver $r \in \{1, 2\}$ and opponent $-r = 3 - r$:
\begin{itemize}
    \item If $\alpha_1 = 0$ and $\alpha_2 = 0$, the round is over: $\Psi_r(E_1, E_2, D, A_{\text{temp}}; \kappa) = R_{\text{safe}}(E_1) - R_{\text{safe}}(E_2)$.
    \item Otherwise, if $A_{\text{temp}} = \emptyset$ or $\alpha_r = 0$, no cards remain for $r$ to execute: $\Psi_r(E_1, E_2, D, A_{\text{temp}}; \kappa) = \kappa(E_1, E_2, D)$.
    \item Otherwise ($A_{\text{temp}} \neq \emptyset$ and $\alpha_r = 1$), let $A' = A_{\text{temp}} \setminus \{c\}$ (multiset difference removing one copy of $c$) and define the nested continuation $\kappa'_c(\tilde{E}_1, \tilde{E}_2, \tilde{D}) = \Psi_r(\tilde{E}_1, \tilde{E}_2, \tilde{D}, A'; \kappa)$. Then
    \[
        \Psi_r(E_1, E_2, D, A_{\text{temp}}; \kappa) = \operatorname{opt}_{r, \, c \in A_{\text{temp}}} \psi_r(E_1, E_2, D, c, \kappa'_c),
    \]
    where:
    \begin{itemize}
        \item \textit{Freeze ($c = \text{fr}$):} target $t^* = -r$ if $\alpha_{-r} = 1$, else $t^* = r$; set $\alpha_{t^*} \leftarrow 0$ and evaluate $\kappa'_c$.
        \item \textit{Second Chance ($c = \text{sc}$):} set $sc_r \leftarrow 1$ if $sc_r = 0$; else set $sc_{-r} \leftarrow 1$ if $\alpha_{-r} = 1$ and $sc_{-r} = 0$; else discard; then evaluate $\kappa'_c$.
        \item \textit{Flip Three ($c = \text{f3}$):} target $t^* = -r$ if $\alpha_{-r} = 1$, else $t^* = r$; evaluate $W_3^{(t^*)}(E_1, E_2, D, \emptyset; \kappa'_c)$.
    \end{itemize}
\end{itemize}
Termination is guaranteed because every recursive call either strictly decreases $|D|$ or preserves $|D|$ while strictly decreasing the total number of deferred action cards across the call stack (bounded above by the $9$ action cards in the deck).

Finally, the solo recursion $W_k^{\text{solo}}(E_2, D, A_{\text{temp}}; \kappa)$ and resolution operator $\Psi^{\text{solo}}(E_2, D, A_{\text{temp}}; \kappa)$ invoked by $V_2$ in Section~2.3.3 are the direct one-player specializations of $W_k^{(2)}$ and $\Psi_2$ (with $\alpha_1 = 0$, Player~2 maximizing $\max_{c \in A_{\text{temp}}}$, busting returning $0$ immediately, and Freeze returning $R_{\text{safe}}(E_2)$ immediately).

\section{The Comprehensive $n$-Player Game}\label{sec:nplayer}
To model the complete round-robin game of Flip 7 for an arbitrary number of players $n \ge 2$, we move beyond a single-agent MDP to a \textbf{finite extensive-form game of perfect information with chance nodes}. Because each player $i \in \{1, \dots, n\}$ optimizes their own strategic utility function $U_i$ over the vector of all players' expected round scores, backward induction from terminal states computes a \textbf{Subgame Perfect Equilibrium (SPE)} (which is unique under the deterministic tie-breaking rule defined below).

\subsection{Card Sets, Joint State Space, and Terminal Scoring}
As in Section~2.3, let $X = \mathcal{N} \cup \mathcal{B} \cup \mathcal{A}_{\text{cards}}$ be the set of all $22$ drawable card types. A global state is a tuple $s = (\mathbf{E}, D)$, where:
\begin{enumerate}
    \item \textbf{Deck Vector $D$:} $D = (d_x)_{x \in X} \in \mathbb{N}_0^{|X|}$ with $|D| = \sum_{x \in X} d_x$ and unit card basis vectors $\boldsymbol{\delta}_x \in \mathbb{N}_0^{|X|}$.
    \item \textbf{Joint Player Environment $\mathbf{E}$:} $\mathbf{E} = (E_1, E_2, \dots, E_n)$, where each player $j \in \{1, \dots, n\}$ has state
    \[
        E_j = (C_j, m_j, p_j, sc_j, \alpha_j)
    \]
    with $C_j \subseteq \mathcal{N}$, $m_j \in \{1, 2\}$, $p_j \in \mathbb{N}_0$, $sc_j \in \{0, 1\}$, and $\alpha_j \in \{0, 1\}$ ($1$ if active, $0$ if stopped, frozen, or busted).
\end{enumerate}

Let $\mathbf{e}_j \in \mathbb{R}^n$ denote the $j$-th standard basis vector in player-payoff space (distinct from the card basis vectors $\boldsymbol{\delta}_x \in \mathbb{N}_0^{|X|}$). Because stopped and frozen players retain their tableaus in $\mathbf{E}$ while busted players are reset to $E^{\text{bust}} = (\emptyset, 1, 0, 0, 0)$ (for which $R_{\text{safe}}(E^{\text{bust}}) = 0$), the complete vector of round scores at any standard terminal state ($\mathcal{T}(\mathbf{E}) = \emptyset$ or $|D| = 0$) is
\[
    \mathbf{R}_{\text{term}}(\mathbf{E}) = \sum_{j=1}^{n} R_{\text{safe}}(E_j)\mathbf{e}_j.
\]
If player $p$ achieves a Flip 7 by drawing a 7th unique number card $x \in \mathcal{N} \setminus C_p$, the round ends immediately with total round-score vector
\[
    \mathbf{R}_{\text{flip7}}(\mathbf{E}, p, x) = R_{\text{flip7}}(E_p, x)\mathbf{e}_p + \sum_{j \neq p} R_{\text{safe}}(E_j)\mathbf{e}_j.
\]
Evaluating all players' scores at terminal states (rather than crediting stayed or frozen players incrementally) ensures that downstream players always observe the full round scores of inactive players a property essential for non-linear utilities such as $U_i^{\text{max}}$ below.

Let $\mathcal{T}(\mathbf{E}) = \{j \in \{1, \dots, n\} \mid \alpha_j = 1\}$ be the set of active players, and let $\mathcal{T}_{-i}(\mathbf{E}) = \mathcal{T}(\mathbf{E}) \setminus \{i\}$. Under the opponent-first targeting rule, the set of valid targets for Freeze or Flip Three played by $i$ is
\[
    \mathcal{T}_{\text{valid}}(i, \mathbf{E}) =
    \begin{cases}
        \mathcal{T}_{-i}(\mathbf{E}) & \text{if } \mathcal{T}_{-i}(\mathbf{E}) \neq \emptyset, \\
        \{i\} & \text{if } \mathcal{T}_{-i}(\mathbf{E}) = \emptyset.
    \end{cases}
\]
(To allow voluntary self-targeting as in official rules, simply replace $\mathcal{T}_{\text{valid}}(i, \mathbf{E})$ with $\mathcal{T}(\mathbf{E})$.) We also write $\mathbf{E}[E_j \leftarrow E_j']$ for the tuple $\mathbf{E}$ with entry $j$ replaced by $E_j'$, and abbreviate single-field updates as $\mathbf{E}[\alpha_j \leftarrow 0]$ or $\mathbf{E}[sc_j \leftarrow 1]$.

\subsection{Utility Functions and the Choice Operator}
Let $\mathbf{V}_i(\mathbf{E}, D) \in \mathbb{R}^n$ denote the vector of \textbf{total expected round scores} for all $n$ players (including scores already banked by inactive players) when it is Player~$i$'s turn in state $(\mathbf{E}, D)$. Player~$i$ evaluates candidate outcome vectors via a utility function $U_i: \mathbb{R}^n \to \mathbb{R}$, such as margin over the average opponent or margin over the leading opponent:
\[
    U_i^{\text{avg}}(\mathbf{v}) = v_i - \frac{1}{n-1} \sum_{j \neq i} v_j, \qquad
    U_i^{\text{max}}(\mathbf{v}) = v_i - \max_{j \neq i} v_j.
\]
\begin{remark}
By linearity of expectation, applying $U_i^{\text{avg}}$ to the expected score vector $\mathbb{E}[\mathbf{S}]$ is exact: $U_i^{\text{avg}}(\mathbb{E}[\mathbf{S}]) = \mathbb{E}[U_i^{\text{avg}}(\mathbf{S})]$. Because $U_i^{\text{max}}$ is non-linear, evaluating $U_i^{\text{max}}(\mathbb{E}[\mathbf{S}])$ serves as a tractable first-moment approximation to expected utility $\mathbb{E}[U_i^{\text{max}}(\mathbf{S})]$, which would otherwise require tracking full joint score distributions rather than expectation vectors.
\end{remark}

Given a finite \textit{ordered tuple} of candidate expected score vectors $\vec{\mathcal{V}} = (\mathbf{v}^{(1)}, \dots, \mathbf{v}^{(m)})$, Player~$i$'s choice operator $\mathcal{C}_i(\vec{\mathcal{V}})$ selects the vector maximizing $U_i$, breaking ties deterministically by smallest index:
\[
    \mathcal{C}_i(\vec{\mathcal{V}}) = \mathbf{v}^{(k^*)}, \qquad
    k^* = \min \Big( \argmax_{k \in \{1, \dots, m\}} U_i(\mathbf{v}^{(k)}) \Big).
\]
Throughout Section~3.3--3.4, candidate tuples follow canonical index orderings: $(\text{stay}, \text{draw})$ for turn actions; ascending player indices $1 < 2 < \dots < n$ for target selection; and $(\text{fr}, \text{sc}, \text{f3})$ for deferred action cards.

\subsection{Round-Robin Turn Order and Backward-Induction Equations}
Play proceeds clockwise among active players. Define $\nextplayer(i, \mathbf{E})$ to return the first element in the cyclic sequence $(i+1, \dots, n, 1, \dots, i)$ belonging to $\mathcal{T}(\mathbf{E})$ (returning $i$ itself if $\mathcal{T}(\mathbf{E}) = \{i\}$), and to return the sentinel $\bot$ if $\mathcal{T}(\mathbf{E}) = \emptyset$.

At any terminal boundary when no active players remain or the deck is empty the round concludes immediately:
\[
    \mathbf{V}_{\bot}(\mathbf{E}, D) = \mathbf{R}_{\text{term}}(\mathbf{E}), \qquad
    \mathbf{V}_i(\mathbf{E}, \mathbf{0}) = \mathbf{R}_{\text{term}}(\mathbf{E}) \quad \forall i \in \{1, \dots, n\}.
\]
When it is active Player~$i$'s turn ($\alpha_i = 1$) with $|D| > 0$, Player~$i$ selects between staying and drawing:
\[
    \mathbf{V}_i(\mathbf{E}, D) = \mathcal{C}_i \Big( \mathbf{Q}_{i,\text{stay}}(\mathbf{E}, D), \, \mathbf{Q}_{i,\text{draw}}(\mathbf{E}, D) \Big).
\]

\begin{itemize}
    \item \textbf{Stay:} Player~$i$ retains their tableau, becomes inactive ($\alpha_i \leftarrow 0$), and passes play to the next active player:
    \[
        \mathbf{Q}_{i,\text{stay}}(\mathbf{E}, D) = \mathbf{V}_{\nextplayer(i, \mathbf{E}')}(\mathbf{E}', D), \qquad \mathbf{E}' = \mathbf{E}[\alpha_i \leftarrow 0].
    \]
    \item \textbf{Draw:}
    \[
        \mathbf{Q}_{i,\text{draw}}(\mathbf{E}, D) = \sum_{\substack{x \in X \\ d_x > 0}} \frac{d_x}{|D|} \cdot \mathbf{v}_{i,\text{next}}(\mathbf{E}, D, x),
    \]
    where $\mathbf{v}_{i,\text{next}}(\mathbf{E}, D, x)$ evaluates the drawn card $x$ as follows:
    \begin{itemize}
        \item \textbf{Duplicate Number ($x \in \mathcal{N} \cap C_i$):}
        \begin{itemize}
            \item \textit{Protected ($sc_i = 1$):} set $\mathbf{E}' = \mathbf{E}[sc_i \leftarrow 0]$ and return $\mathbf{V}_{\nextplayer(i, \mathbf{E}')}(\mathbf{E}', D - \boldsymbol{\delta}_x)$.
            \item \textit{Bust ($sc_i = 0$):} set $\mathbf{E}' = \mathbf{E}[E_i \leftarrow E^{\text{bust}}]$ and return $\mathbf{V}_{\nextplayer(i, \mathbf{E}')}(\mathbf{E}', D - \boldsymbol{\delta}_x)$.
        \end{itemize}

        \item \textbf{Safe Number ($x \in \mathcal{N} \setminus C_i$):}
        \begin{itemize}
            \item \textit{Flip 7 ($|C_i| = 6$):} the round ends immediately with $\mathbf{v}_{i,\text{next}} = \mathbf{R}_{\text{flip7}}(\mathbf{E}, i, x)$.
            \item \textit{Standard Safe Draw ($|C_i| < 6$):} set $\mathbf{E}' = \mathbf{E}[C_i \leftarrow C_i \cup \{x\}]$ and return $\mathbf{V}_{\nextplayer(i, \mathbf{E}')}(\mathbf{E}', D - \boldsymbol{\delta}_x)$.
        \end{itemize}

        \item \textbf{Bonus Card ($x \in \mathcal{B}$):} set $\mathbf{E}' = \mathbf{E}[m_i \leftarrow 2m_i]$ (if $x \in \mathcal{B}_{\times}$) or $\mathbf{E}' = \mathbf{E}[p_i \leftarrow p_i + \val(x)]$ (if $x \in \mathcal{B}_{+}$), and return $\mathbf{V}_{\nextplayer(i, \mathbf{E}')}(\mathbf{E}', D - \boldsymbol{\delta}_x)$.

        \item \textbf{Freeze ($x = \text{fr}$):} Player~$i$ selects a target $t \in \mathcal{T}_{\text{valid}}(i, \mathbf{E})$ (ordered by ascending $t$) to deactivate while keeping $E_t$ intact:
        \[
            \mathbf{v}_{i,\text{next}} = \mathcal{C}_i \Big( \big( \mathbf{V}_{\nextplayer(i, \mathbf{E}[\alpha_t \leftarrow 0])}(\mathbf{E}[\alpha_t \leftarrow 0], D - \boldsymbol{\delta}_x) \big)_{t \in \mathcal{T}_{\text{valid}}(i, \mathbf{E})} \Big).
        \]

        \item \textbf{Second Chance ($x = \text{sc}$):} Let $\mathcal{T}_{\text{sc}}(i, \mathbf{E}) = \{j \in \mathcal{T}_{-i}(\mathbf{E}) \mid sc_j = 0\}$. Then
        \[
            \mathbf{v}_{i,\text{next}} =
            \begin{cases}
                \mathbf{V}_{\nextplayer(i, \mathbf{E}[sc_i \leftarrow 1])}(\mathbf{E}[sc_i \leftarrow 1], D - \boldsymbol{\delta}_x) & \text{if } sc_i = 0, \\
                \mathcal{C}_i \Big( \big( \mathbf{V}_{\nextplayer(i, \mathbf{E}[sc_t \leftarrow 1])}(\mathbf{E}[sc_t \leftarrow 1], D - \boldsymbol{\delta}_x) \big)_{t \in \mathcal{T}_{\text{sc}}(i, \mathbf{E})} \Big) & \text{if } sc_i = 1,\, \mathcal{T}_{\text{sc}}(i, \mathbf{E}) \neq \emptyset, \\
                \mathbf{V}_{\nextplayer(i, \mathbf{E})}(\mathbf{E}, D - \boldsymbol{\delta}_x) & \text{if } sc_i = 1,\, \mathcal{T}_{\text{sc}}(i, \mathbf{E}) = \emptyset.
            \end{cases}
        \]

        \item \textbf{Flip Three ($x = \text{f3}$):} Define the macro-turn continuation $\kappa_i(\mathbf{E}', D') = \mathbf{V}_{\nextplayer(i, \mathbf{E}')}(\mathbf{E}', D')$. Player~$i$ chooses an active target $t \in \mathcal{T}_{\text{valid}}(i, \mathbf{E})$:
        \[
            \mathbf{v}_{i,\text{next}} = \mathcal{C}_i \Big( \big( \mathbf{W}_3^{(t)}(\mathbf{E}, D - \boldsymbol{\delta}_x, \emptyset; \, \kappa_i) \big)_{t \in \mathcal{T}_{\text{valid}}(i, \mathbf{E})} \Big).
        \]
    \end{itemize}
\end{itemize}

\subsection{Generalizing the Flip Three Sub-Game and Action-Card Resolution}
As in Section~2.4, we define the vector-valued forced-draw recursion $\mathbf{W}_k^{(t)}(\mathbf{E}, D, A_{\text{temp}}; \kappa) \in \mathbb{R}^n$ and resolution operator $\mathbf{\Psi}_r(\mathbf{E}, D, A_{\text{temp}}; \kappa) \in \mathbb{R}^n$ parameterized by a continuation $\kappa(\mathbf{E}', D') \in \mathbb{R}^n$.

\subsubsection{The Recursive Forced-Draw Step ($\mathbf{W}_k^{(t)}$)}
When $k = 0$ or the deck runs dry ($|D| = 0$), forced draws conclude and target $t$ resolves $A_{\text{temp}}$:
\[
    \mathbf{W}_k^{(t)}(\mathbf{E}, D, A_{\text{temp}}; \kappa) = \mathbf{\Psi}_t(\mathbf{E}, D, A_{\text{temp}}; \kappa) \qquad \text{if } k = 0 \text{ or } |D| = 0.
\]
For $k \in \{1, 2, 3\}$ and $|D| > 0$:
\[
    \mathbf{W}_k^{(t)}(\mathbf{E}, D, A_{\text{temp}}; \kappa) = \sum_{\substack{x \in X \\ d_x > 0}} \frac{d_x}{|D|} \cdot \mathbf{w}_{\text{step}}^{(t)}(x),
\]
where:
\begin{itemize}
    \item \textbf{Action Card ($x \in \mathcal{A}_{\text{cards}}$):} $\mathbf{w}_{\text{step}}^{(t)}(x) = \mathbf{W}_{k-1}^{(t)}(\mathbf{E}, D - \boldsymbol{\delta}_x, A_{\text{temp}} \cup \{x\}; \kappa)$.
    \item \textbf{Bonus Card ($x \in \mathcal{B}$):} update $\mathbf{E}' = \mathbf{E}[m_t \leftarrow 2m_t]$ (if $x \in \mathcal{B}_{\times}$) or $\mathbf{E}' = \mathbf{E}[p_t \leftarrow p_t + \val(x)]$ (if $x \in \mathcal{B}_{+}$), and set $\mathbf{w}_{\text{step}}^{(t)}(x) = \mathbf{W}_{k-1}^{(t)}(\mathbf{E}', D - \boldsymbol{\delta}_x, A_{\text{temp}}; \kappa)$.
    \item \textbf{Safe Number ($x \in \mathcal{N} \setminus C_t$):}
    \[
        \mathbf{w}_{\text{step}}^{(t)}(x) =
        \begin{cases}
            \mathbf{R}_{\text{flip7}}(\mathbf{E}, t, x) & \text{if } |C_t| = 6, \\
            \mathbf{W}_{k-1}^{(t)}\big(\mathbf{E}[C_t \leftarrow C_t \cup \{x\}], D - \boldsymbol{\delta}_x, A_{\text{temp}}; \, \kappa\big) & \text{if } |C_t| < 6.
        \end{cases}
    \]
    \item \textbf{Duplicate Number ($x \in \mathcal{N} \cap C_t$):}
    \begin{itemize}
        \item \textit{Protected ($sc_t = 1$):} $\mathbf{w}_{\text{step}}^{(t)}(x) = \mathbf{W}_{k-1}^{(t)}(\mathbf{E}[sc_t \leftarrow 0], D - \boldsymbol{\delta}_x, A_{\text{temp}}; \kappa)$.
        \item \textit{Bust \& Transfer ($sc_t = 0$):} Let $\mathbf{E}_{\text{bust}} = \mathbf{E}[E_t \leftarrow E^{\text{bust}}]$. If $\mathcal{T}(\mathbf{E}_{\text{bust}}) = \emptyset$, the round ends immediately with $\mathbf{R}_{\text{term}}(\mathbf{E}_{\text{bust}})$. If $A_{\text{temp}} = \emptyset$ and $\mathcal{T}(\mathbf{E}_{\text{bust}}) \neq \emptyset$, control passes directly to $\kappa(\mathbf{E}_{\text{bust}}, D - \boldsymbol{\delta}_x)$. Otherwise, busted Player~$t$ chooses an active recipient $j \in \mathcal{T}(\mathbf{E}_{\text{bust}})$ via $\mathcal{C}_t$ to inherit $A_{\text{temp}}$:
        \[
            \mathbf{w}_{\text{step}}^{(t)}(x) = \mathcal{C}_t \Big( \big( \mathbf{\Psi}_j(\mathbf{E}_{\text{bust}}, D - \boldsymbol{\delta}_x, A_{\text{temp}}; \, \kappa) \big)_{j \in \mathcal{T}(\mathbf{E}_{\text{bust}})} \Big).
        \]
    \end{itemize}
\end{itemize}

\subsubsection{The Action-Card Resolution Operator ($\mathbf{\Psi}_r$)}
For an acting resolver $r \in \{1, \dots, n\}$:
\begin{itemize}
    \item If $\mathcal{T}(\mathbf{E}) = \emptyset$, the round ends immediately: $\mathbf{\Psi}_r(\mathbf{E}, D, A_{\text{temp}}; \kappa) = \mathbf{R}_{\text{term}}(\mathbf{E})$.
    \item Otherwise, if $A_{\text{temp}} = \emptyset$ or $\alpha_r = 0$, return $\kappa(\mathbf{E}, D)$.
    \item Otherwise ($A_{\text{temp}} \neq \emptyset$ and $\alpha_r = 1$), Player~$r$ chooses which card $c \in A_{\text{temp}}$ to execute next (ordered canonically as $\text{fr}, \text{sc}, \text{f3}$). Letting $A' = A_{\text{temp}} \setminus \{c\}$ and $\kappa'_c(\tilde{\mathbf{E}}, \tilde{D}) = \mathbf{\Psi}_r(\tilde{\mathbf{E}}, \tilde{D}, A'; \kappa)$:
    \[
        \mathbf{\Psi}_r(\mathbf{E}, D, A_{\text{temp}}; \kappa) = \mathcal{C}_r \Big( \big( \boldsymbol{\psi}_r(\mathbf{E}, D, c, \kappa'_c) \big)_{c \in A_{\text{temp}}} \Big),
    \]
    where each single-card resolution $\boldsymbol{\psi}_r$ mirrors the macro-turn rules:
    \begin{itemize}
        \item \textit{Freeze ($c = \text{fr}$):} $\mathcal{C}_r \big( ( \kappa'_c(\mathbf{E}[\alpha_t \leftarrow 0], D) )_{t \in \mathcal{T}_{\text{valid}}(r, \mathbf{E})} \big)$.
        \item \textit{Second Chance ($c = \text{sc}$):} $\kappa'_c(\mathbf{E}[sc_r \leftarrow 1], D)$ if $sc_r = 0$; else $\mathcal{C}_r \big( ( \kappa'_c(\mathbf{E}[sc_t \leftarrow 1], D) )_{t \in \mathcal{T}_{\text{sc}}(r, \mathbf{E})} \big)$ if $\mathcal{T}_{\text{sc}}(r, \mathbf{E}) \neq \emptyset$; else $\kappa'_c(\mathbf{E}, D)$.
        \item \textit{Flip Three ($c = \text{f3}$):} $\mathcal{C}_r \big( ( \mathbf{W}_3^{(t)}(\mathbf{E}, D, \emptyset; \kappa'_c) )_{t \in \mathcal{T}_{\text{valid}}(r, \mathbf{E})} \big)$.
    \end{itemize}
\end{itemize}
Because every recursive transition either strictly decreases $|D|$ or strictly decreases the total count of pending action cards, the recursion is strictly acyclic and uniquely defines the SPE value vectors $\mathbf{V}_i(\mathbf{E}, D)$ for all states.



\section{Computational Tractability and Approximations}
While backward induction guarantees an exact optimal policy (Section~1--2) and a unique Subgame Perfect Equilibrium (Section~\ref{sec:nplayer}), the practical feasibility of exact dynamic programming depends heavily on whether the draw pile is tracked from a fixed initial deck or allowed to vary arbitrarily across rounds, as well as on the number of players $n$.

\subsection{State-Space Complexity}
Across all formulations, the primary bottleneck is the dimensionality of the deck vector $D \in \mathbb{N}_0^{|X|}$ combined with the exponential growth of the joint player environment $\mathbf{E} = (E_1, \dots, E_n)$.

\begin{enumerate}
    \item \textbf{Unconstrained Deck Space:} If the draw pile carries over across multiple rounds without reshuffling, the number of distinct deck vectors is the product of $(d_{x,0} + 1)$ over all card types $x \in X$:
    \begin{itemize}
        \item \textit{Number cards only ($79$ cards):} $2 \cdot \prod_{x=1}^{12}(x + 1) = 2 \cdot 13! \approx 1.25 \times 10^{10}$ deck states.
        \item \textit{Full Flip 7 deck ($94$ cards):} $2 \cdot 13! \cdot 2^{6} \cdot 4^{3} = 8{,}192 \cdot 13! \approx 5.10 \times 10^{13}$ deck states.
    \end{itemize}
    Tabulating value functions over the entire unconstrained deck space is therefore infeasible in memory.

    \item \textbf{Single-Player Game from a Fixed Starting Deck ($D_0$):} In a single round starting from a known deck $D_0$, every safe draw adds a unique card to the player's tableau and removes it from $D_0$, so the remaining deck $D$ is uniquely determined by the player's drawn cards:
    \begin{itemize}
        \item \textit{Simplified (number cards only):} A non-terminal hand $C \subseteq \mathcal{N}$ satisfies $|C| \le 6$, yielding exactly $\sum_{k=0}^{6} \binom{13}{k} = 2^{13}/2 = 4{,}096$ reachable states.
        \item \textit{With bonus cards:} Tracking which subset of the $6$ bonus cards has been drawn adds at most a factor of $2^6 = 64$, giving at most $4{,}096 \times 64 = 262{,}144$ reachable states.
    \end{itemize}
    Both single-player models can be solved exactly in milliseconds on standard hardware.

    \item \textbf{Two-Player Sequential Game from a Fixed Deck ($D_0$):} In the simplified number-only model (Section~2.1), Player~1 leaves behind either a ``stay'' deck (at most $4{,}096$ distinct decks, each missing up to $6$ unique cards) or a ``bust'' deck (missing $1$ to $6$ unique cards plus a second copy of one card $x \in \{2, \dots, 12\}$, yielding $11 \sum_{j=0}^{5} \binom{12}{j} = 17{,}446$ disjoint decks). Thus Player~2 inherits at most $21{,}542$ distinct deck states. Because each sub-problem $V_2(\emptyset, D)$ explores at most $4{,}096$ hands, the entire simplified two-player game requires at most $21{,}542 \times 4{,}096 \approx 8.82 \times 10^{7}$ state evaluations and negligible memory if each $V_2(\emptyset, D)$ is solved independently making it fully tractable. Once bonus and action cards are introduced, however, Second Chance discards and nested Flip Three draws decouple $D$ from $(E_1, E_2)$, inflating the reachable state space by several orders of magnitude.

    \item \textbf{Comprehensive $n$-Player Game:} Each individual player environment $E_j = (C_j, m_j, p_j, sc_j, \alpha_j)$ admits at most $4{,}096$ non-terminal sets $C_j$, $2$ multipliers, $16$ achievable flat-bonus sums ($p_j \in \{0, 2, \dots, 30\}$), $2$ Second Chance states, and $2$ activity states up to $|\mathcal{E}| \approx 5.24 \times 10^{5}$ states per player. Because the joint state $(\mathbf{E}, D)$ scales as $\mathcal{O}(|\mathcal{E}|^{n} \cdot |\mathcal{D}_{\text{reach}}|)$ (even after enforcing card-conservation constraints), exact backward induction suffers from the curse of dimensionality and becomes computationally intractable for $n \ge 3$.
\end{enumerate}

\begin{table}[h]
\centering
\begin{tabular}{l l l}
\hline
\textbf{Model Variant (from fixed $D_0$)} & \textbf{Reachable States / Evaluations} & \textbf{Exact Solvability} \\
\hline
1-Player, number cards only & $4{,}096$ states & Trivial ($< 1\text{ ms}$) \\
1-Player, with bonus cards & $\le 2.62 \times 10^{5}$ states & Fast ($< 1\text{ s}$) \\
2-Player sequential, number cards only & $\le 8.82 \times 10^{7}$ evaluations & Tractable (seconds) \\
2-Player sequential, full deck & $\gg 10^{10}$ states & Intractable without pruning \\
$n$-Player round-robin ($n \ge 3$), full deck & $\mathcal{O}(|\mathcal{E}|^{n} |\mathcal{D}_{\text{reach}}|)$ & Intractable \\
\hline
\end{tabular}
\caption{Summary of computational complexity for exact backward induction across Flip 7 model variants.}
\label{tab:complexity}
\end{table}

\subsection{Practical Approximation Methods}
To obtain strong policies for the full two-player and $n$-player games where exact backward induction is infeasible, four complementary approximation strategies are well suited to Flip 7:

\begin{enumerate}
    \item \textbf{Deck Abstraction (Sampling-with-Replacement and Feature Compression):} Rather than tracking the full $22$-dimensional vector $D$, one can approximate card-draw probabilities using either the static initial distribution (ignoring depletion) or a low-dimensional deck summary such as remaining deck size $|D|$, current bust probability $\sum_{x \in C_i} d_x / |D|$, and remaining counts of action cards. This eliminates the $\sim 5.10 \times 10^{13}$ deck factor while retaining the dominant risk signal.

    \item \textbf{Decoupled Opponent Modeling (Mean-Field Approximation):} For $n \ge 3$, tracking the full tuple $\mathbf{E}$ can be replaced by conditioning Player~$i$'s policy on their own environment $E_i$ plus compact opponent summary statistics: the highest currently secured or live opponent score $\max_{j \neq i} R_{\text{safe}}(E_j)$, the count of remaining active opponents $|\mathcal{T}_{-i}(\mathbf{E})|$, and baseline expected scores for active opponents computed from the exact single-player model.

    \item \textbf{Online Tree Search (Expectimax and Monte Carlo Tree Search):} Because a single round of Flip 7 has short horizon depth (typically $10$--$30$ draws across all players), decisions can be computed online from the live state $(\mathbf{E}, D)$ without precomputing the full state space. Chance-node Monte Carlo Tree Search (MCTS) or depth-limited Expectimax using the fast single-player value function $V(C_j, m_j, p_j, D)$ as a leaf-node heuristic directly handles action-card targeting and Flip Three chains.

    \item \textbf{Approximate Dynamic Programming and Self-Play:} The equilibrium value vector $\mathbf{V}_i(\mathbf{E}, D)$ can be parameterized by a function approximator $\widehat{\mathbf{V}}_i(\mathbf{E}, D; \theta)$ (e.g., a neural network or linear combination of domain features such as bust risk, distance to Flip 7 $7 - |C_i|$, multiplier status, and score gap to the leader) and trained via multi-agent self-play reinforcement learning or fitted value iteration.
\end{enumerate}

\printbibliography

\end{document}